\documentclass[12pt,a4paper]{article}

\usepackage[utf8]{inputenc}
\usepackage[T1]{fontenc}
\usepackage{geometry}
\usepackage{setspace}
\usepackage{titlesec}
\usepackage{enumitem}
\usepackage{booktabs}
\usepackage{hyperref}

\geometry{margin=1in}
\onehalfspacing

\title{\textbf{Activities and Functions of the Statistical System in India and Kerala}}
\author{Gayathri P S}
\date{Sep,2026}

\begin{document}

\maketitle

\section{Introduction}

The official statistical system in India operates at both central and state levels. At the national level, the Ministry of Statistics and Programme Implementation (MoSPI), Government of India, serves as the principal authority for official statistics. At the state level, the Department of Economics and Statistics (DES), Government of Kerala, functions as the nodal agency for official statistics in the state. Both institutions collect, compile, analyse and disseminate statistical information for planning, policy formulation, monitoring and evidence-based decision-making.

\section{Ministry of Statistics and Programme Implementation, Government of India}

MoSPI coordinates the development of India's national statistical system and produces major economic and socio-economic indicators. Through the National Statistical Office (NSO), it compiles National Accounts such as Gross Domestic Product (GDP), Gross Value Added (GVA), national income, consumption, savings and capital formation. It also produces important indicators such as the Consumer Price Index (CPI) and Index of Industrial Production (IIP).

MoSPI conducts large-scale surveys including the Periodic Labour Force Survey (PLFS), Household Consumption Expenditure Survey (HCES), Comprehensive Annual Modular Survey (CAMS), Multiple Indicator Survey (MIS) and Annual Survey of Industries (ASI). These surveys provide information on employment, household consumption, industrial activity and living conditions. The Ministry also develops statistical methodologies, sampling and estimation techniques, standards and classifications, and disseminates datasets and statistical publications.

Through its Programme Implementation Wing, MoSPI monitors Central Sector infrastructure projects costing more than Rs.~150 crore and oversees the Members of Parliament Local Area Development Scheme (MPLADS).

\section{Department of Economics and Statistics, Government of Kerala}

DES Kerala is responsible for producing state-level and local-level statistics covering population, agriculture, industry, prices and socio-economic development. Under Population and Demographics, it compiles Vital Statistics on births and deaths and publishes Medical Certification of Cause of Death (MCCD) statistics.

In agriculture, DES functions as the State Agricultural Statistics Authority (SASA). It undertakes the Establishment of an Agency for Reporting Agricultural Statistics (EARAS), Agriculture Census, Cost of Cultivation studies and evaluation of soil conservation. These activities provide information on crop area, production, yield, agricultural holdings and cultivation costs.

For industrial statistics, DES conducts or coordinates the Annual Survey of Industries (ASI) and compiles the Index of Industrial Production (IIP). It estimates State Domestic Product (SDP/GSDP), district income and per-capita income. The department also collects daily prices of essential commodities and produces Consumer Price Index data, price bulletins and market reviews.

Its socio-economic activities include National Sample Surveys, Labour and Housing Statistics, Economic Census, ad-hoc surveys, infrastructure statistics, environment statistics and gender statistics. DES disseminates information through Handbooks, Economic Reviews, Fact Sheets, Local Level Statistics, project reports and newsletters, along with digital statistical databases and dashboards.

\section{Conclusion}

MoSPI and DES Kerala perform complementary roles in India's statistical system. MoSPI provides national-level economic and socio-economic indicators and coordinates the broader statistical framework, while DES Kerala generates detailed state, district and local-level statistics. Through surveys, censuses, economic estimates, price monitoring and statistical publications, both institutions provide reliable evidence for planning, policy formulation, resource allocation, programme monitoring and evaluation.

\end{document}
